\documentclass[11pt]{article}
\usepackage[a4paper,margin=1in]{geometry}
\usepackage{amsmath,amssymb,mathtools,bm}
\usepackage{enumitem,booktabs,array}
\usepackage{tcolorbox}
\usepackage{hyperref}
\usepackage[protrusion=true,expansion=false]{microtype}
\hypersetup{colorlinks=true,linkcolor=blue,urlcolor=blue,citecolor=blue}
\setlist[itemize]{topsep=2pt,itemsep=2pt}
\setlist[enumerate]{topsep=2pt,itemsep=2pt}
\newtcolorbox{keybox}{colback=blue!3!white,colframe=blue!40!black,boxrule=0.4pt,arc=1mm}
\newtcolorbox{alertbox}{colback=red!3!white,colframe=red!40!black,boxrule=0.4pt,arc=1mm}
\newtcolorbox{answerbox}{colback=green!3!white,colframe=green!40!black,boxrule=0.4pt,arc=1mm}
\newtcolorbox{drillbox}{colback=yellow!5!white,colframe=orange!60!black,boxrule=0.4pt,arc=1mm}
\newcommand{\EV}{\mathbb{E}}
\newcommand{\Var}{\mathrm{Var}}
\newcommand{\Cov}{\mathrm{Cov}}
\newcommand{\blank}[1]{\underline{\hspace{#1}}}

\title{E01-04: Zingales (2000, JF) \\
\large ``In Search of New Foundations'' --- Active Recall Workbook}
\author{Empirical Corporate Finance II --- Exam Prep}
\date{\today}
\begin{document}\maketitle

\begin{keybox}
\textbf{Paper in one sentence.} Classical theories of the firm (Coase, transaction costs, property rights) assume the firm's key assets are physical and alienable, but modern firms increasingly rely on \emph{human capital} and \emph{relationships} --- non-alienable critical resources --- so corporate-finance foundations must be rebuilt around who has access to and control over these critical non-owned resources.

\textbf{Placement.} AFA Presidential Address. Sets the agenda for 2000s+ corporate-finance theory: human capital, employee ownership, networks, innovation. Paired with Rajan-Zingales (1998, 2001) on access/power as firm foundation.
\end{keybox}

\section*{Round 1: Complete Walk-Through}

\subsection*{1.1 The puzzle}
Traditional (Grossman-Hart-Moore) theory: the firm is a collection of physical assets; residual control over those assets determines boundaries, incentives, and value. But:
\begin{itemize}
\item Service, software, and knowledge firms have few physical assets.
\item Key value is created by employees with mobile human capital.
\item Markets increasingly rely on long-term relationships, not spot contracts.
\end{itemize}
A property-rights view on physical assets is less and less descriptive.

\subsection*{1.2 Proposed foundation: critical resources and access}
\begin{itemize}
\item Define the firm by its \emph{critical resources} --- assets the firm cannot easily replace or substitute.
\item These may include human capital (key employees), brand, relationships, data, information.
\item \emph{Access} to a critical resource (the ability to work with it) generates power similar to, but distinct from, ownership.
\item Firms and market power flow from \emph{control of access}, not necessarily title ownership.
\end{itemize}

\subsection*{1.3 Implications for corporate-finance questions}
\begin{itemize}
\item \textbf{Financing.} Debt/equity decisions must account for the firm's dependence on non-owned human capital; equity-like rewards (stock options) substitute for cash wages.
\item \textbf{Ownership structure.} Why do firms give equity to employees? To bond them to the firm (access $\to$ ownership).
\item \textbf{Boundaries.} The firm may outsource physical production but retain critical relationships (brand, design).
\item \textbf{Governance.} The board and shareholders protect only part of the surplus; labor and other stakeholders claim the rest via access.
\end{itemize}

\subsection*{1.4 Relation to earlier work}
\begin{itemize}
\item Rajan-Zingales (1998): \emph{Power in a Theory of the Firm} -- access and specific knowledge.
\item Rajan-Zingales (2001): \emph{Influence Costs and Hierarchy} -- implications for internal organization.
\item Contrast with Grossman-Hart: shifts from owning physical assets to controlling critical resources (including human).
\end{itemize}

\subsection*{1.5 Agenda and predictions}
\begin{itemize}
\item More human-capital-intensive firms should have flatter hierarchies and more employee equity.
\item Brand- and relationship-intensive firms should integrate just enough to retain access.
\item Debt levels should be lower where key resources are human capital (non-liquidatable).
\item Innovation and software firms are natural test beds.
\end{itemize}

\subsection*{1.6 Caveats}
\begin{alertbox}
\begin{itemize}
\item Verbal framework, not a formal model; specific predictions are inherited from RZ 1998/2001.
\item ``Critical resources'' definition is elastic --- hard to measure empirically.
\item Does not fully replace GH: physical-asset ownership still matters for capital-intensive sectors.
\end{itemize}
\end{alertbox}

\section*{Round 2: Level 1 Fill-in}
\begin{enumerate}
\item Traditional theory assumes assets are \blank{2cm} and alienable.
\item Modern firms rely on \blank{2cm} capital and relationships.
\item Zingales proposes rebuilding around \blank{2cm} resources and \blank{1.5cm} to them.
\item Employee equity bonds workers to firm-\blank{2cm} access.
\item Debt levels should be \blank{1cm} when key resources are human capital.
\end{enumerate}

\section*{Round 3: Level 2 Prompts}
\begin{enumerate}
\item State Zingales's puzzle: why does GH property-rights theory under-describe modern firms?
\item Define critical resources and access; contrast with ownership.
\item Discuss financing and boundary implications of the new foundation.
\item Relate Zingales (2000) to Rajan-Zingales (1998, 2001).
\item Propose an empirical test contrasting human-capital vs.\ physical-capital firms.
\end{enumerate}

\section*{Round 4: Minimal Scaffolding}
From a blank page: (i) puzzle, (ii) critical resources + access, (iii) financing implications, (iv) RZ relation, (v) agenda.

\section*{Round 5: Blank Page}
Essay: new foundations for the theory of the firm in a human-capital economy.

\vspace{6pt}
\begin{drillbox}
\textbf{Speed Drill.} (1) Puzzle. (2) Proposed core unit. (3) Ownership vs.\ access. (4) One financing prediction. (5) Companion paper.
\end{drillbox}

\section*{Quick Reference \& Answer Key}
\begin{answerbox}
\begin{itemize}
\item \textbf{Puzzle.} Modern firms have few alienable physical assets.
\item \textbf{Core unit.} Critical resources + access.
\item \textbf{Implication.} Employee equity, lower debt, brand-centric boundaries.
\item \textbf{Paired.} Rajan-Zingales 1998, 2001.
\end{itemize}
\end{answerbox}

\begin{keybox}
\textbf{30-second exam pitch.} Zingales (2000), his AFA presidential address, argues the Grossman-Hart-Moore property-rights foundation --- where the firm is a set of owned physical assets and residual control confers value --- under-describes modern firms, whose key assets are human capital, brand, relationships, and data, which are non-alienable. He proposes rebuilding the foundation around \emph{critical resources} (assets the firm cannot easily substitute for) and \emph{access} to them, where access generates power akin to but distinct from ownership. Corporate-finance implications: employee equity bonds human capital to the firm; debt is lower where critical resources are human; boundaries include relationship/brand not just physical plant; governance must accommodate claimants whose power flows from access rather than equity ownership. Paired with Rajan-Zingales (1998, 2001), it sets the 2000s+ research agenda on firms built on human capital, innovation, and networks, and motivates empirical work contrasting human-capital-intensive and physical-capital-intensive firms.
\end{keybox}

\end{document}
